% Options for packages loaded elsewhere
\PassOptionsToPackage{unicode}{hyperref}
\PassOptionsToPackage{hyphens}{url}
\documentclass[
]{article}
\usepackage{xcolor}
\usepackage[margin=1in]{geometry}
\usepackage{amsmath,amssymb}
\setcounter{secnumdepth}{-\maxdimen} % remove section numbering
\usepackage{iftex}
\ifPDFTeX
  \usepackage[T1]{fontenc}
  \usepackage[utf8]{inputenc}
  \usepackage{textcomp} % provide euro and other symbols
\else % if luatex or xetex
  \usepackage{unicode-math} % this also loads fontspec
  \defaultfontfeatures{Scale=MatchLowercase}
  \defaultfontfeatures[\rmfamily]{Ligatures=TeX,Scale=1}
\fi
\usepackage{lmodern}
\ifPDFTeX\else
  % xetex/luatex font selection
\fi
% Use upquote if available, for straight quotes in verbatim environments
\IfFileExists{upquote.sty}{\usepackage{upquote}}{}
\IfFileExists{microtype.sty}{% use microtype if available
  \usepackage[]{microtype}
  \UseMicrotypeSet[protrusion]{basicmath} % disable protrusion for tt fonts
}{}
\makeatletter
\@ifundefined{KOMAClassName}{% if non-KOMA class
  \IfFileExists{parskip.sty}{%
    \usepackage{parskip}
  }{% else
    \setlength{\parindent}{0pt}
    \setlength{\parskip}{6pt plus 2pt minus 1pt}}
}{% if KOMA class
  \KOMAoptions{parskip=half}}
\makeatother
\usepackage{longtable,booktabs,array}
\newcounter{none} % for unnumbered tables
\usepackage{calc} % for calculating minipage widths
% Correct order of tables after \paragraph or \subparagraph
\usepackage{etoolbox}
\makeatletter
\patchcmd\longtable{\par}{\if@noskipsec\mbox{}\fi\par}{}{}
\makeatother
% Allow footnotes in longtable head/foot
\IfFileExists{footnotehyper.sty}{\usepackage{footnotehyper}}{\usepackage{footnote}}
\makesavenoteenv{longtable}
\usepackage{graphicx}
\makeatletter
\newsavebox\pandoc@box
\newcommand*\pandocbounded[1]{% scales image to fit in text height/width
  \sbox\pandoc@box{#1}%
  \Gscale@div\@tempa{\textheight}{\dimexpr\ht\pandoc@box+\dp\pandoc@box\relax}%
  \Gscale@div\@tempb{\linewidth}{\wd\pandoc@box}%
  \ifdim\@tempb\p@<\@tempa\p@\let\@tempa\@tempb\fi% select the smaller of both
  \ifdim\@tempa\p@<\p@\scalebox{\@tempa}{\usebox\pandoc@box}%
  \else\usebox{\pandoc@box}%
  \fi%
}
% Set default figure placement to htbp
\def\fps@figure{htbp}
\makeatother
\setlength{\emergencystretch}{3em} % prevent overfull lines
\providecommand{\tightlist}{%
  \setlength{\itemsep}{0pt}\setlength{\parskip}{0pt}}
\usepackage{bookmark}
\IfFileExists{xurl.sty}{\usepackage{xurl}}{} % add URL line breaks if available
\urlstyle{same}
% fallback for those not using the hyperref driver hyperxmp:
\makeatletter
\@ifundefined{xmpquote}{\newcommand{\xmpquote}[1]{#1}}{}
\makeatother
\hypersetup{
  pdftitle={Statistics Anxiety and Academic Attainment},
  pdfauthor={Emma Crawford, Alice Miller, Leah Cookson, Alexandra Anderson},
  hidelinks,
  pdfcreator={LaTeX via pandoc}}

\title{Statistics Anxiety and Academic Attainment}
\author{Emma Crawford, Alice Miller, Leah Cookson, Alexandra Anderson}
\date{2026-11-19}

\begin{document}
\maketitle

\subsection{Effect Size Summary}\label{effect-size-summary}

\emph{to be completed using our final analysis:}

\(d = A, \ 95\%\ \mathrm{CI}\ [B,C]\), based on \(N = D\).

\subsection{Summary of the Analysis
Approach}\label{summary-of-the-analysis-approach}

\emph{to be completed together after choosing the primary analysis (max
100 words)}

\subsection{Details of the Analysis}\label{details-of-the-analysis}

\subsubsection{Data Preparation and Participant
Exclusions}\label{data-preparation-and-participant-exclusions}

\paragraph{Initial Inspection and Grade
Availability}\label{initial-inspection-and-grade-availability}

The SMARVUS dataset contained 12,570 observations and 788 variables. We
checked for identical rows and missing or duplicated participant
identifiers (`unique\_id'); none were found. Missing reponses were
investigated with reference to the variables required for this analysis
rather than excluding participants with missing entries in other places
of the questionnaire. Statistics attainment was recorded in `grade',
with 1820 participants having a grade entry other than `NA' or `\,``Not
Available''\,'.

\paragraph{Standarising Attainment Across
Institutions}\label{standarising-attainment-across-institutions}

Reported grades were recorded differently across the participating
surveys. This included different numerical scales, letter grades and
ordinal categories. We examined each institutions lowest and highest
documented grades before we contrusted a comparable attainment variable.
Supported numerical marks (0--100, 0--20, 0--10 and 1--10) were linearly
rescaled to a common 0--100 range:

\[
\text{Scaled grade} = 100\,\frac{\text{Observed grade} - \text{Lowest possible grade}}{\text{Highest possible grade} - \text{Lowest possible grade}}.
\]

To keep observations that used letter grades, we usedthe midpoint of the
applicable institution-specific percentage band as an \emph{estimated}
numerical mark. Grade categories that could not be converted reasonably
or defensibly, including the 1--5 ordinal systems, ambiguous entries and
values ourside continuous-attainment datasets. This left us with 1362
converted numerical grades and 320 estimated letter grades (1662
observations in total). Grade-band midpoints are approximations, and it
is important to note than converting to a common range does not remove
the differences in grading standards across different institutions.

\paragraph{Attention Checks and Exclusion
Criteria}\label{attention-checks-and-exclusion-criteria}

We assessed response quality using the six attention check items and the
`attention\_amnesty' response. Participants were excluded if they
answered at least one administered check incorrectly or explicitly
selected `\,``No''\,' to the amnesty question. Missing check or amnesty
responses were \textbf{not} interpreted as failures as the questions
were not necessarily recorded for every survey. Under these rules, 1542
participants remained after attainment and quality screening. This is
the \emph{data-preparation sample}, not necessarily the final analysis
sample: the anxiety-scale completeness criteria will be applied
subsequently.

The participant counts at each stage are summarised in the table below.
Of the 1542 kept observations, 165 had no recorded attention check
answers and 70 had a missing amnesty response. So, the conclusions
should be checked for the sensitivity to the treatment of missing
quality information. Results could also depend on if the estimated
letter grades are included.

{\def\LTcaptype{none} % do not increment counter
\begin{longtable}[]{@{}lrr@{}}
\toprule\noalign{}
Stage & N & Excluded \\
\midrule\noalign{}
\endhead
\bottomrule\noalign{}
\endlastfoot
Original dataset & 12570 & NA \\
Non-missing grade entries & 1820 & 10750 \\
Usable standardised attainment & 1682 & 138 \\
After attention amnesty & 1655 & 27 \\
After attention checks & 1542 & 113 \\
\end{longtable}
}

\subsubsection{Anxiety measures and
scoring}\label{anxiety-measures-and-scoring}

\emph{To be completed by Alice.}

\subsubsection{Statistical analysis and effect
size}\label{statistical-analysis-and-effect-size}

\emph{To be completed by Leah.}

\subsubsection{Mathematics anxiety
comparison}\label{mathematics-anxiety-comparison}

\emph{To be completed by Alex.}

\subsubsection{Interpretation and
limitations}\label{interpretation-and-limitations}

\emph{To be completed together.}

RMARS is revised mathematics anxiety rating scale (NORMAL MATHS ANXIETY
ONE) Q8.2\_1 to Q8.2\_24 RMARS-S is like RMARS but replaces the maths
with stats Q7.2\_1 to Q7.2\_23 STARS is statistics anxiety rating scale
(NORMAL STATS ANXIETY ONE) Q7.1\_1 to Q7.1\_23 STARS-M is a modified
version of STARS but replaces the stats with maths Q8.1\_1 to Q8.1\_20

\subsection{Effect size summary}\label{effect-size-summary-1}

Simply state:

d = A, 95\% CI {[}B, C{]} based on N = D.

where A, B, C are replaced with the values that you have found using
data from D participants.

\subsection{Summary of the analysis
approach}\label{summary-of-the-analysis-approach-1}

Note that the \texttt{echo\ =\ FALSE} parameter was added to the code
chunk to prevent printing of the R code that generated the plot.

Give a short summary of your approach, specifying the type of analysis
you used and any other important choices you made (e.g., about which
data to include/exclude).

\subsection{Details of the analysis}\label{details-of-the-analysis-1}

This is the main body of the report, which should make clear how you
arrived at the values quoted in the first section. This should be in a
similar style to the lecture notes for this course, with a mix of text,
R code, and output generated by the code. It may be helpful to think of
the intended audience as a classmate who has not yet engaged with the
project -- you can assume a good level of statistical background
knowledge, but you will need to give a concise explanation of the
important features of the dataset and your specific analysis choices.
There is no specific page limit, since it will depend on the amount of
code and plots included, but should be unlikely to need more than 6
pages.

\end{document}
